% ભાગ: first-order-logic
% પ્રકરણ: models-theories
% વિભાગ: expressing-props-of-structures

\documentclass[../../../include/open-logic-section]{subfiles}

\begin{document}

\olfileid{fol}{mat}{exs}

\olsection{\printtoken{P}{structure}ના ગુણધર્મો વ્યક્ત કરવા}

\begin{explain}
વિધેયો અને સંબંધો પરની શરતો, અથવા વધુ સામાન્ય રીતે કોઈ સંરચનામાંનાં
વિધેયો અને સંબંધો આ શરતોને સંતોષે છે એ વ્યક્ત કરવું ઘણી વાર ઉપયોગી અને
અગત્યનું હોય છે. દાખલા તરીકે, કોઈ ભાષા માટેની જે !!{structure}sમાં
!!{predicate}sનો આપણને જોઈતો ``અર્થ'' મળે છે તેમને એવી
સંરચનાઓથી અલગ પાડવાની રીતો આપણને જોઈએ છે જેમાં તે અર્થ મળતો
નથી. અલબત્ત, કઈ !!{structure}sને આપણે ``અભિપ્રેત'' ગણીએ તે કહેવાની
આપણને સંપૂર્ણ છૂટ છે. જેમ કે, !!{predicate}~$\le$નું અર્થઘટન ક્રમ જ
હોવું જોઈએ એમ કહી શકીએ, અથવા આપણને~$\Lang L$નાં માત્ર એવાં અર્થઘટનોમાં
રસ છે જેમાં વ્યાપકક્ષેત્ર ગણોનું બનેલું હોય અને~$\Obj \in$નું
અર્થઘટન ``નો !!a{element} છે'' સંબંધ તરીકે થાય. પરંતુ શું આ વાત ભાષાનાં
!!{sentence}s વડે કરી શકીએ? બીજા શબ્દોમાં, !!a{structure}~$\Struct M$
પરની કઈ શરતોને~$\Struct M$ની ભાષાના !!a{sentence} (અથવા કદાચ
!!{sentence}sના ગણ) વડે વ્યક્ત કરી શકીએ? કેટલીક શરતો આપણે વ્યક્ત કરી
શકીશું નહિ. દાખલા તરીકે,~$\Lang L_A$નું એવું કોઈ વાક્ય નથી જે
!!{structure}~$\Struct M$માં માત્ર $\Domain M=\Nat$ હોય ત્યારે જ સત્ય
હોય. ``વ્યાપકક્ષેત્રમાં માત્ર પ્રાકૃતિક સંખ્યાઓ જ છે'' એ આપણે વ્યક્ત
કરી શકતા નથી. છતાં !!{structure}sના કેટલાક ``સંરચનાત્મક ગુણધર્મો''
કદાચ વ્યક્ત કરી શકીએ. !!{structure}sના કયા ગુણધર્મોને !!{sentence}s
વડે વ્યક્ત કરી શકીએ? અથવા બીજી રીતે કહીએ તો, !!{structure}sના કયા
સંગ્રહને કોઈ !!a{sentence} (અથવા !!{sentence}sના ગણ)ને સત્ય બનાવતી
સંરચનાઓ તરીકે વર્ણવી શકીએ?
\end{explain}

\begin{defn}[ગણનું નિદર્શ]
$\Gamma$, ભાષા~$\Lang L$નાં !!{sentence}sનો ગણ હોય. જો દરેક
$!A\in\Gamma$ માટે $\Sat{M}{!A}$, તો આપણે કહીએ કે
!!a{structure}~$\Struct M$,~$\Gamma$નું \emph{નિદર્શ છે}.
\end{defn}

\begin{ex}
વાક્ય~$\lforall[x][x \le x]$,~$\Struct M$માં ત્યારે અને માત્ર ત્યારે
સત્ય છે જ્યારે $\Assign{\le}{M}$ સ્વવાચક સંબંધ હોય. વાક્ય
$\lforall[x][\lforall[y][((x \le y \land y \le x) \lif x = y)]]$,
~$\Struct M$માં ત્યારે અને માત્ર ત્યારે સત્ય છે જ્યારે
$\Assign{\le}{M}$ વિસંમિત હોય. વાક્ય
$\lforall[x][\lforall[y][\lforall[z][((x \le y \land y \le z)
      \lif x \le z)]]]$,~$\Struct M$માં ત્યારે અને માત્ર ત્યારે સત્ય
છે જ્યારે $\Assign{\le}{M}$ પરંપરિત હોય. તેથી
\begin{align*}
\{\quad &\lforall[x][x \le x], \\
   & \lforall[x][\lforall[y][((x \le y \land y \le
    x) \lif x = y)]], \\
   &\lforall[x][\lforall[y][\lforall[z][((x \le y
      \land y \le z) \lif x \le z)]]] \quad \}
\end{align*}
નાં નિદર્શો બરાબર એવી સંરચનાઓ છે જેમાં~$\Assign{\le}{M}$ સ્વવાચક,
વિસંમિત અને પરંપરિત, એટલે કે આંશિક ક્રમ, હોય. તેથી તેમને
\emph{આંશિક ક્રમોના પ્રથમ-ક્રમ સિદ્ધાંત}નાં સ્વયંસિદ્ધિઓ તરીકે લઈ
શકીએ છીએ.
\end{ex}

\end{document}
